\documentclass[11pt]{article}
\usepackage[utf8]{inputenc}
\usepackage{amsmath,amssymb}
\usepackage{hyperref}
\title{Efficient Microbiome Host Interaction: A Resource-Aware Approach}
\author{Anonymous}
\date{}
\begin{document}
\maketitle

\begin{abstract}
This paper studies Microbiome Host Interaction. Grounded in a real, citable literature \cite{ref1,ref2,ref3,ref4,ref5,ref6,ref7,ref8},
we propose a focused method and report \textbf{illustrative} experiments.
All references below are real and verifiable (OpenAlex/arXiv); the reported
numbers are simulated to illustrate intended behaviour.
\end{abstract}

\section{Introduction}
Microbiome Host Interaction is an active research area. This work targets a concrete gap identified from
the recent literature and contributes a method together with an evaluation protocol.

\section{Related Work (real literature)}
The following works, retrieved from public bibliographic APIs, frame our study:
\begin{itemize}
\item Berendsen et al. (2012) — "The rhizosphere microbiome and plant health", Trends in Plant Science (cited 5390).
\item Zheng et al. (2020) — "Interaction between microbiota and immunity in health and disease", Cell Research (cited 3943).
\item Trivedi et al. (2020) — "Plant–microbiome interactions: from community assembly to plant health", Nature Reviews Microbiology (cited 3790).
\item Lundberg et al. (2012) — "Defining the core Arabidopsis thaliana root microbiome", Nature (cited 2963).
\item Coyte et al. (2015) — "The ecology of the microbiome: Networks, competition, and stability", Science (cited 2884).
\item Desai et al. (2016) — "A Dietary Fiber-Deprived Gut Microbiota Degrades the Colonic Mucus Barrier and Enhances Pathogen Susceptibility", Cell (cited 2880).
\end{itemize}

\section{Method}
We reduce the dominant computational cost via adaptive resource allocation, activating only the components needed per input. We instantiate this idea for Microbiome Host Interaction and analyse its cost/quality trade-off.

\section{Experiments (Illustrative)}
\begin{center}
\begin{tabular}{lcc}
System & Primary Metric & Efficiency \\ \hline
Strong baseline & 0.812 & 1.00$\times$ \\
Ablation & 0.828 & 0.92$\times$ \\
\textbf{Ours} & \textbf{0.851} & \textbf{0.71$\times$}
\end{tabular}
\end{center}
\emph{Metrics simulated to illustrate the intended effect; not measured.}

\section{Conclusion}
We connected a real literature to a concrete method for Microbiome Host Interaction. Future work will
validate the illustrative results with real experiments.

\begin{thebibliography}{9}
\bibitem{ref1} Roeland L. Berendsen, Corné M. J. Pieterse, Peter A. H. M. Bakker. The rhizosphere microbiome and plant health. \emph{Trends in Plant Science}, 2012. DOI:10.1016/j.tplants.2012.04.001
\bibitem{ref2} Danping Zheng, Timur Liwinski, Eran Elinav. Interaction between microbiota and immunity in health and disease. \emph{Cell Research}, 2020. DOI:10.1038/s41422-020-0332-7
\bibitem{ref3} Pankaj Trivedi, Jan E. Leach, Susannah G. Tringe et al.. Plant–microbiome interactions: from community assembly to plant health. \emph{Nature Reviews Microbiology}, 2020. DOI:10.1038/s41579-020-0412-1
\bibitem{ref4} Derek S. Lundberg, Sarah L. Lebeis, Sur Herrera Paredes et al.. Defining the core Arabidopsis thaliana root microbiome. \emph{Nature}, 2012. DOI:10.1038/nature11237
\bibitem{ref5} Katharine Z. Coyte, Jonas Schlüter, Kevin R. Foster. The ecology of the microbiome: Networks, competition, and stability. \emph{Science}, 2015. DOI:10.1126/science.aad2602
\bibitem{ref6} Mahesh S. Desai, Anna M. Seekatz, Nicole M. Koropatkin et al.. A Dietary Fiber-Deprived Gut Microbiota Degrades the Colonic Mucus Barrier and Enhances Pathogen Susceptibility. \emph{Cell}, 2016. DOI:10.1016/j.cell.2016.10.043
\bibitem{ref7} Mary E. Rinella, Brent A. Neuschwander‐Tetri, Mohammad Shadab Siddiqui et al.. AASLD Practice Guidance on the clinical assessment and management of nonalcoholic fatty liver disease. \emph{Hepatology}, 2023. DOI:10.1097/hep.0000000000000323
\bibitem{ref8} Jonathan J. Havel, Diego Chowell, Timothy A. Chan. The evolving landscape of biomarkers for checkpoint inhibitor immunotherapy. \emph{Nature reviews. Cancer}, 2019. DOI:10.1038/s41568-019-0116-x
\end{thebibliography}
\end{document}
